\section{Methods and Experimental Setup}
\label{sec:methods}

\subsection{Problem formulation}
We study learned conditional channel simulation.
Let \(x \in \mathbb{R}^n\) denote the transmitted channel input and \(y \in \mathbb{R}^n\) the corresponding channel output.
We consider two target parameterizations and two experimental protocols.
The direct-output parameterization learns \(p(y\mid x)\) and predicts \(\hat y=g_\theta(x,z)\).
The residual parameterization defines
\begin{equation}
e = y - x.
\end{equation}
In residual mode, generative models approximate \(p(e \mid x)\), and channel outputs are reconstructed as
\begin{equation}
\hat{y} = x + \hat{e}.
\end{equation}
The generator-level protocol compares the proposed drifting family against conditional DDPM, DDIM, and WGAN references on AWGN, Rayleigh, and SSPA, where matched reference runs are available.
The drift-field ablation then holds the direct-output drifting architecture fixed and compares target-kernel, joint-kernel, joint Sinkhorn, and condition-wise Sinkhorn training fields over 100 seeds, except for the SSPA update-budget check described below.
That ablation is evaluated with direct SWD, anchor-conditioned diagnostics, and downstream symbolic SER/BER, and additionally includes a compact TDL wireless channel.
The distinction between direct-output and residual targets matters because the two protocols define different modeling biases: residual learning can simplify near-identity channels, while direct-output learning aligns with the conditional law used by diffusion, WGAN, and the downstream channel implants.

\subsection{Conditional diffusion baseline}
The diffusion baseline follows the channel-learning setting of \cite{paper2309} and uses a conditional DDPM objective \cite{ddpm} with DDIM sampling for deterministic inference trajectories \cite{ddim}.
The backbone is a conditional multilayer perceptron defined by
\begin{equation}
\nu_\theta(y_t, t, x),
\end{equation}
where \(y_t\) is the noised channel output at diffusion step \(t\), and \(x\) is the channel condition.

The network architecture is a fully connected model with three hidden conditional layers and one linear output head.
In the benchmark implementation, the first layer takes the concatenated vector \([y_t, x] \in \mathbb{R}^{2n}\).
Each hidden layer is modulated by a learned embedding of the diffusion step \(t\), and a softplus nonlinearity is applied after each conditional affine transform.
The final layer maps the hidden representation back to \(\mathbb{R}^n\).

Training uses a cosine noise schedule with horizon \(T=100\) in the publication-grade configuration.
The model is trained in \(v\)-prediction mode:
\begin{equation}
\mathcal{L}_{\text{diff}}
=
\mathbb{E}_{y_0, x, t, \epsilon}
\left[
\left\|
\nu - \nu_\theta(y_t, t, x)
\right\|_2^2
\right].
\end{equation}
For AWGN, the reference diffusion run uses \(30\) epochs at learning rate \(10^{-4}\).
For Rayleigh, it uses a two-stage schedule of \(10\) epochs at \(10^{-3}\) followed by \(20\) epochs at \(10^{-4}\).
For SSPA, the reference setup uses a longer training budget.
At inference time, we report both the stochastic DDPM reverse chain and the deterministic DDIM trajectory.
The DDIM trajectory uses skip-based index sets such as
\( \{9,19,\ldots,99\} \) for DDIM(10) and
\( \{4,9,\ldots,99\} \) for DDIM(20).
The main comparison in the final figures uses DDPM(100) and DDIM(100); reduced-step DDIM is retained as an auxiliary ablation.

\subsection{Conditional drifting baseline}
The drifting baseline is a conditional one-shot generator inspired by \cite{drifting}.
Following the notation of the drifting paper, let \(f_\theta\) map \(z \sim p_z\) to a sample \(u = f_\theta(z)\), and let
\begin{equation}
q_\theta = (f_\theta)_\# p_z
\end{equation}
denote the induced pushforward distribution.
The central object is a drifting field \(V_{p,q}(u)\) that indicates how samples from the current model distribution \(q\) should move toward the target distribution \(p\).
The intended equilibrium property is that the field vanishes when the two distributions match:
\begin{equation}
q = p
\quad \Longrightarrow \quad
V_{p,q}(u) = 0.
\end{equation}
This viewpoint leads to a fixed-point style training update.
Starting from \(u = f_\theta(z)\), a drifted target is formed as
\begin{equation}
\tilde{u}
=
\mathrm{stopgrad}\!\left(
u + \eta V_{p,q_\theta}(u)
\right),
\end{equation}
and the generator is optimized by
\begin{equation}
\mathcal{L}_{\text{drift}}
=
\mathbb{E}_{z \sim p_z}
\left[
\left\|
f_\theta(z) - \tilde{u}
\right\|_2^2
\right].
\end{equation}
This makes clear why inference is naturally one-shot: the iterative correction is absorbed into training, and sampling after training requires only one evaluation of \(f_\theta\) \cite{drifting}.

For channel simulation, two target parameterizations are meaningful.
The first is the direct conditional channel-output law \(p(y \mid x)\), which treats the learned generator as a surrogate for the full channel law.
The second is the residual law \(p(e \mid x)\) with \(e = y - x\), which injects the inductive bias that the channel output can be understood as the transmitted symbol plus a structured perturbation.
The direct formulation is conceptually cleaner and corresponds to learning the channel law itself.
The residual formulation is usually easier to optimize on near-identity channels, because the model does not need to relearn the identity component of the channel map.

In direct mode, the generator predicts \(\hat{y} = g_\theta(x, z)\) and is evaluated against samples from \(p(y \mid x)\).
In residual mode, the generator predicts \(\hat{e} = g_\theta(x, z)\) and the channel output is reconstructed as \(\hat{y} = x + \hat{e}\).
In the experiments, the direct formulation is used when comparing against direct-output diffusion and WGAN baselines and throughout the Sinkhorn drift-field ablation.
The residual formulation is retained as a complementary parameterization because it exposes when the choice of target space changes the result.

In the residual benchmark instantiation, the implemented generator is therefore
\begin{equation}
g_\theta(x, z),
\qquad
z \sim \mathcal{N}(0, I),
\end{equation}
and predicts the residual \(\hat{e} = g_\theta(x, z)\).
The generator \(g_\theta\) is implemented as a two-hidden-layer multilayer perceptron with hidden width \(128\), latent dimension \(16\), and SiLU activations:
\begin{equation}
[x, z]
\rightarrow
\mathrm{Linear}
\rightarrow
\mathrm{SiLU}
\rightarrow
\mathrm{Linear}
\rightarrow
\mathrm{SiLU}
\rightarrow
\mathrm{Linear}.
\end{equation}

The training target is constructed from an attraction--repulsion drift field on the chosen target representation.
This is our concrete channel-specific instantiation of the generic field \(V_{p,q}\) from the drifting formulation.
Given generated samples \(G = \{g_i\}\) and positive samples \(P = \{p_j\}\) in either direct-output space or residual space, we compute
\begin{equation}
V(g_i)
=
\left(
\sum_j w^{(P)}_{ij} p_j - g_i
\right)
-
\lambda
\left(
\sum_k w^{(G)}_{ik} g_k - g_i
\right),
\end{equation}
where \(w^{(P)}\) and \(w^{(G)}\) are row-normalized RBF kernel weights.
The first term attracts generated samples toward nearby positive samples, while the second term repels samples away from local generator collapse.
The training target is detached:
\begin{equation}
\tilde{g}_i
=
\mathrm{stopgrad}\!\left(g_i + \eta V(g_i)\right),
\end{equation}
and the generator is trained by
\begin{equation}
\mathcal{L}_{\text{drift}}
=
\frac{1}{N}
\sum_i
\left\|
g_i - \tilde{g}_i
\right\|_2^2.
\end{equation}
Thus, the generic update \(u \mapsto u + \eta V_{p,q}(u)\) becomes a kernel-based update estimated directly from minibatches of generated and true channel samples in the selected target space.

We also include an enhanced direct-output variant with a conditioning-aware kernel.
For this variant, sample interactions are computed on joint condition-target pairs \(k((x_i,y_i),(x_j,y_j))\) rather than on \(k(y_i,y_j)\) alone.
This makes attraction and repulsion depend on both transmitted-symbol similarity and channel-output similarity, and the direct-output target kernel can additionally include residual similarity.
The generator architecture and one-shot inference path are unchanged.

We additionally consider a Sinkhorn replacement for this heuristic kernel field, motivated by W-Flow \cite{han2026wflow}.
For channel simulation, the Sinkhorn field is applied conditionally rather than on the unrestricted joint space.
Writing \(p(dx,dy)=\mu(dx)p_x(dy)\) and \(q_\theta(dx,dy)=\mu(dx)q_{\theta,x}(dy)\), the conditional objective is
\begin{equation}
\mathcal S_\varepsilon^{\mathrm{cond}}(q_\theta,p)
=
\int S_\varepsilon(q_{\theta,x},p_x)\,\mu(dx).
\end{equation}
The resulting condition-wise velocity is
\begin{equation}
v_x(y)
=
T^\varepsilon_{q_{\theta,x},p_x}(y)
-
T^\varepsilon_{q_{\theta,x},q_{\theta,x}}(y),
\end{equation}
where the barycentric projections are computed from entropic optimal-transport couplings.
For simulator channels, this is estimated by drawing multiple true and generated samples at the same condition anchor.
For measured data with only one output per condition, the corresponding approximation is local conditional Sinkhorn over nearby condition blocks.

Empirically, the two parameterizations behave differently.
At matched mid-budget settings, residual drifting is materially stronger than direct \(y \mid x\) drifting on AWGN, Rayleigh, and SSPA, which is consistent with the view that residual modeling supplies a useful near-identity inductive bias.
For this reason, both parameterizations are reported: direct drifting is used for direct-output comparisons and the Sinkhorn ablation, while residual drifting is retained as a diagnostic for near-identity channel structure.

The corrected benchmark uses the following hyperparameters:
\(\eta = 1.0\),
repulsive weight \(\lambda = 1.0\),
minimum kernel bandwidth \(0.2\),
and maximum drift norm \(2.0\).
Unlike diffusion, inference requires only one forward pass through \(g_\theta(x, z)\), so the per-sample inference cost does not scale with \(T\).

\subsection{Conditional GAN baseline}
Two adversarial baselines are used.
For the direct-output benchmark, we include a conditional WGAN baseline matching the comparison setup of \cite{paper2309}.
For the residual diagnostic benchmark, we include a stronger conditional GAN adapted from the nonlinear-channel code path.

In this benchmark, the generator receives the concatenated vector \([x, z]\), where \(z \sim \mathcal{N}(0, I)\), and outputs the full channel output \(\hat{y}\).
Architecturally, the generator uses an input linear layer, two batch-normalized hidden blocks, a central hidden layer, two additional batch-normalized hidden blocks, and a residual skip from the condition input \(x\) to the output.
The discriminator takes the concatenated pair \([y, x]\), processes it with a residual-style multilayer perceptron, concatenates the original input back into the hidden state, and produces a scalar realism score.

Two adversarial objectives are implemented in the benchmark code: WGAN-GP and a stronger feature-aligned GAN mode denoted \texttt{GAN\_FA}.
The residual diagnostic benchmark uses \texttt{GAN\_FA}, since it was substantially more stable than the plain WGAN-GP baseline in this channel-learning setting.
In \texttt{GAN\_FA}, the discriminator is trained with BCE-with-logits and label smoothing, and the generator is trained with an adversarial term plus an \(L_1\) reconstruction term:
\begin{equation}
\mathcal{L}_{G}
=
\mathcal{L}_{\text{adv}}
+
\alpha
\mathbb{E}\left[\|\hat{y} - y\|_1\right].
\end{equation}
In the experiments, \(\alpha = 1.0\), generator updates are delayed relative to discriminator updates, and the default hidden width is \(128\).

\subsection{Channel models}
We benchmark four channel families.
All channels are implemented on real-valued vectors, but not all channels use the same interpretation of the coordinates.
For AWGN and Rayleigh, the benchmark uses real-valued channel inputs and outputs directly.
In the main benchmark, both channels use dimension \(n=7\).
For SSPA and OptFib, the signal is again stored as a real vector, but adjacent entries are interpreted as paired in-phase and quadrature components.
Thus SSPA uses \(n=8\), corresponding to four I/Q pairs, and OptFib uses \(n=2\), corresponding to one I/Q pair.
Noise variables are sampled independently across real components unless the channel transformation itself introduces coupling.

\textbf{AWGN:}
\begin{equation}
y = x + n,
\qquad
n \sim \mathcal{N}(0,\sigma^2 I).
\end{equation}
This is the standard additive white Gaussian noise reference channel and serves as the simplest near-identity baseline.
In our benchmark, AWGN is treated as a real-valued vector channel rather than as a complex I/Q channel.

\textbf{Rayleigh:}
the transmitted symbol is multiplied elementwise by a random fading amplitude and then corrupted by additive Gaussian noise.
In the implemented benchmark,
\begin{equation}
y = h \odot x + n,
\end{equation}
where \(\odot\) denotes elementwise multiplication,
\begin{equation}
h_k = \frac{1}{\sqrt{2}}\sqrt{u_k^2 + v_k^2},
\qquad
u_k,v_k \overset{\text{i.i.d.}}{\sim} \mathcal{N}(0,1),
\end{equation}
and $n \sim \mathcal{N}(0,\sigma^2 I).$

Each real component is therefore scaled by an independent Rayleigh-distributed magnitude before noise is added.
This channel is multiplicative rather than additive and is correspondingly less aligned with a pure residual perturbation view.
As implemented here, Rayleigh is also a real-valued vector channel rather than a paired-I/Q complex channel.

\textbf{SSPA:}
the input is passed through a smooth solid-state power amplifier nonlinearity before additive noise is applied.
In this case, the real-valued representation is interpreted as paired I/Q components.
Let
\begin{equation}
r = \sqrt{x_I^2 + x_Q^2}
\end{equation}
denote the input amplitude.
The nonlinear amplitude law used in the benchmark is
\begin{equation}
a(r)
=
\frac{G}{\left(1+\left(\frac{Gr}{A_0}\right)^{2p}\right)^{1/(2p)}},
\end{equation}
with default parameters \(p=3\), \(A_0=1.5\), and \(G=5.0\).
The noiseless amplifier output is $\tilde{x} = a(r)\,x,$
and the observed channel output is
\begin{equation}
y = \tilde{x} + n,
\qquad
n \sim \mathcal{N}\!\left(0,\frac{\sigma^2}{2}I\right).
\end{equation}
This channel introduces amplitude-dependent nonlinear distortion while preserving the input direction in the I/Q plane.

\textbf{OptFib:}
We include a simplified memoryless optical-fiber proxy channel following \cite{li2018e2efiber} for the choice of parameters.
As in SSPA, the real-valued state is interpreted as an I/Q pair.
Let the state after propagation step \(k\) be \(x^{(k)} = [x_I^{(k)},x_Q^{(k)}]^\top\).
At each step, the instantaneous nonlinear phase is
\begin{equation}
\phi^{(k)}
=
\gamma
\left(\left(x_I^{(k)}\right)^2 + \left(x_Q^{(k)}\right)^2\right)\frac{L}{K_{\text{step}}}.
\end{equation}
The deterministic Kerr-style rotation is
\begin{equation}
\begin{aligned}
\tilde{x}_I^{(k+1)} &= x_I^{(k)}\cos\phi^{(k)} - x_Q^{(k)}\sin\phi^{(k)}, \\
\tilde{x}_Q^{(k+1)} &= x_I^{(k)}\sin\phi^{(k)} + x_Q^{(k)}\cos\phi^{(k)},
\end{aligned}
\end{equation}
followed by additive Gaussian perturbations at every step:
\begin{equation}
\begin{aligned}
x_I^{(k+1)} &= \tilde{x}_I^{(k+1)} + \sigma_n \xi_I^{(k)}, \\
x_Q^{(k+1)} &= \tilde{x}_Q^{(k+1)} + \sigma_n \xi_Q^{(k)},
\end{aligned}
\qquad
\xi_I^{(k)},\xi_Q^{(k)} \overset{\text{i.i.d.}}{\sim} \mathcal{N}(0,1).
\end{equation}
When that optical-channel parameterization is used,
\begin{equation}
\sigma_n
=
\sqrt{\frac{10^{(P_n-30)/10}}{2K_{\text{step}}}},
\end{equation}
with \(P_n\) measured in dBm.


\subsection{Training protocol}
We use two complementary protocols.
First, the generator-level protocol compares the proposed drifting family against conditional DDPM, DDIM, and WGAN references on AWGN, Rayleigh, and SSPA.
Second, the drift-field ablation holds the direct-output drifting architecture fixed and compares target-space kernel drifting, conditioning-aware joint-kernel drifting, joint Sinkhorn drifting, and condition-wise Sinkhorn drifting over 100 seeds.
The ablation reports direct SWD, anchor-conditioned diagnostics, and downstream symbolic SER/BER.
The TDL channel is included in this second protocol because it tests a memory-bearing wireless channel, while no matched DDPM/DDIM/WGAN TDL reference run is reported.

For AWGN and Rayleigh, the setting is \(n=7\); for SSPA and TDL, \(n=8\).
The downstream symbolic autoencoder uses message-alphabet size \(M_{\mathrm{msg}}=16\) for AWGN, Rayleigh, and TDL, and \(M_{\mathrm{msg}}=64\) for SSPA.
Thus \(R=\log_2(M_{\mathrm{msg}})/n\), giving rates \(4/7\), \(4/7\), \(6/8\), and \(4/8\), respectively.
The \((E_b/N_0,R)\) pairs are \((5\,\mathrm{dB},4/7)\) for AWGN, \((12\,\mathrm{dB},4/7)\) for Rayleigh, \((8\,\mathrm{dB},6/8)\) for SSPA, and \((10\,\mathrm{dB},4/8)\) for TDL.

For the generator-level AWGN and Rayleigh runs, the training dataset size is \(10^7\), the batch size is \(5000\), and drifting, diffusion, and WGAN are trained for \(30\) epochs.
For SSPA, the dataset size is \(10^7\), the batch size is \(4096\), and the training budget is \(160\) epochs.
The reported SSPA condition-wise Sinkhorn row uses the same epoch count but a compact \(120{,}000\)-sample budget, which corresponds to approximately \(4.8\)k optimizer updates and matches the observed convergence of the corrected same-condition field.
TDL uses the compact nonlinear-channel budget: dataset size \(120{,}000\), batch size \(512\), \(60\) drifting epochs, and evaluation size \(100{,}000\).
Downstream symbolic autoencoders are trained through either the analytic channel or the learned channel implant and are always evaluated on the analytic channel.

\subsection{Evaluation metric}
The primary generator metric is sliced Wasserstein distance (SWD), computed in the target space used by each benchmark path.
For direct-output models, we report SWD on generated and true channel outputs.
For residual models, we report SWD on generated and true residual vectors:
\begin{equation}
\mathrm{SWD}\!\left(\{\hat{e}_i\}, \{e_i\}\right).
\end{equation}
Given two empirical sample clouds \(P=\{e_i\}_{i=1}^N\) and \(Q=\{\hat e_i\}_{i=1}^N\), SWD is computed by drawing random unit directions \(\theta_m \in \mathbb{S}^{n-1}\), projecting both clouds onto each line, and averaging the one-dimensional Wasserstein distances:
\begin{equation}
\mathrm{SWD}(P,Q)
=
\frac{1}{M}
\sum_{m=1}^M
W_1\!\left(\{\langle \theta_m,e_i\rangle\}_{i=1}^N,\{\langle \theta_m,\hat e_i\rangle\}_{i=1}^N\right).
\end{equation}
In other words, SWD measures how different the generated and true sample clouds look when viewed from many random one-dimensional directions.
This makes it sensitive to both location and shape discrepancies while remaining much cheaper to estimate than a full high-dimensional Wasserstein distance.
In the generator-level runs, SWD is estimated with \(128\) random projections for AWGN, Rayleigh, and SSPA.
The drift-field ablation additionally reports anchor-conditioned SWD and Gaussian-W2 diagnostics, where repeated outputs are sampled at fixed transmitted symbols.
These conditional metrics are reported because global output-space SWD can disagree with downstream symbolic SER/BER.

\subsection{Inference-time comparison}
The latency comparison is measured separately from the fidelity benchmark.
For the stored CPU timing run, DDPM(100) required \(4.2612\) seconds per repeat over 20,480 samples, while the one-shot drifting generator required \(0.0163\) seconds under the same batch structure.
This corresponds to an observed speedup of approximately \(260.8\times\) in favor of drifting.

\subsection{Reproducibility}
The benchmark script exposes an explicit random seed argument and the publication runner defaults to deterministic execution unless disabled.
Additional sanity checks include metric identity and symmetry checks for SWD, byte-identical reruns under fixed seed, finite-value validation of stored JSON result files, and small cross-seed ranking checks.
These checks are not a substitute for the full publication-scale run, but they provide confidence that the benchmark pipeline is internally consistent.
